\subsection*{1. 定义 (Definition)}
对某个数学概念或术语的精确、清晰、无歧义的描述。它通过列出一个事物必须具备的所有且仅有的属性来界定其含义。定义是理论的基础，解决了概念的范畴问题。
\begin{itemize}
    \item 一般句式为：“满足 XX 条件的 B 称为 A。”
\end{itemize}

\subsection*{2. 公理 (Axiom / Postulate)}
在理论体系中，一个未经证明、但被当作不证自明的基本命题。其真实性被视为理所当然，并作为演绎及推论其他事实的起点。例如：“过两个不重合的点有且仅有一条直线。”
\begin{itemize}
    \item 公理是理论的基石，使得理论能够被人的理性所接受。
\end{itemize}

\subsection*{3. 定理 (Theorem)}
一个在公理和定义的基础上，被严格的数学逻辑证明为真的重要陈述。在数学研究中，“定理”通常用于指代那些最重要或核心的结论。证明定理是数学的中心活动。

\subsection*{4. 命题 (Proposition)}
一个可以判断真假的陈述句。在数学语境下，它通常指一个已被证明为真的结论，但其重要性或普适性不及“定理”。定理和命题的区别主要在于理论高度和重要性。

\subsection*{5. 引理 (Lemma)}
一个本身可能不太重要，但其唯一目的是为了辅助证明一个更重要定理的命题。当一个定理的证明过程比较复杂或冗长时，常将其中一些关键的中间步骤独立出来作为引理先行证明，从而使主定理的证明结构更清晰。
\begin{itemize}
    \item 引理是证明复杂定理的“铺路石”或“基石”。
\end{itemize}

\subsection*{6. 推论 (Corollary)}
一个能够从某个定理或命题中通过简单、直接的推导得出的结论。推论通常紧随其所依赖的定理之后出现，是该定理的直接应用或自然延伸。

\subsection*{7. 猜想 (Conjecture)}
一个被相信为真，但尚未被证明的数学陈述。许多猜想（如哥德巴赫猜想）是数学研究的重要驱动力。一旦一个猜想被证明，它就成为了一个定理。

\subsection*{8. 断言 (Claim)}
一个在证明过程中提出的断定，并随即给出其证明。它通常像一个非正式的、局部的引理，用于澄清证明中的某个具体步骤。

\subsection*{9. 证明 (Proof)}
使用公理、定义以及已被证明的定理，通过有效的逻辑论证，来解释为什么一个陈述（如定理、命题）为真的过程。

\subsection*{10. 恒等式 (Identity)}
一个表示两个（通常含变量的）数学表达式恒等的方程。例如三角恒等式、欧拉恒等式等。

\subsection*{11. 悖论 (Paradox)}
指在一个给定的公理体系下，可以被同时证明为真又为假的陈述。悖论的存在往往揭示了理论体系中可能存在的缺陷或不一致性（如罗素悖论）。该词有时也用于描述那些与直觉相悖但逻辑上成立的惊人结论。

\subsection*{12. 假说 (Hypothesis)}
根据已知的科学事实和原理，对所研究的现象及其规律性提出的推测性说明。它是有待验证的理论起点，比“猜想”更侧重于探索和解释。

\subsection*{13. 记法 (Notation)}
用于表示数学概念、对象、运算和逻辑关系的符号体系。良好的记法是数学表达清晰、简洁的基础。